\chapter{Computers: Four-Loop, Fourth-Order Systems (4-OS)}\label{lect6}
\localtoc

After the first part of this textbook described the symbolic structures involved in computing, and the second part described the physical structures that can be used to implement computing systems, the time has come to involve the unifying entity represented by {\bf {\em information}}. Indeed, information is the one that takes control once the serendipitous meeting between representations, physical structures, and an acting symbolic structure occurs. The symbolic structure we are referring to is the information that manifests itself in the form of computer programs. Computing systems represent the synthesis between the physical structure of circuits and the symbolic structure of programs. This was made possible by a conceptual clarification -- the emergence of mathematical models of computation -- and a technological evolution -- the structuring of digital systems by increasing their autonomy through the impact of included loops.


\section{The mathematical model of the Turing Machine}

In current language we use the term computation without a rigorous definition. When we have to substantiate the concept of computability and computation, rigor is absolutely necessary. We will continue with not only a historical but also a theoretical introduction to the concept of computation in the sense in which it is conceived in computational science.

\subsection{The Decision Problem}

Despite appearances, computers are not numerical instruments but logical structures that have their theoretical origin in solving decision-making problems. To the question: {\em where and when does the history of computers begin?}, the answer, however strange it may seem, is: {\em somewhere on the island of Crete about two and a half millennia ago}. Indeed, Epimenides of Knossos (the capital of the island of Crete), a semi-mythical Greek seer from the 7th or 6th century BC, caught himself saying that {\em "all Cretans are liars"}. Surprised, because in the next second he could not {\bf decide} whether the assertion he had just made was true or false. Thus was born, if the story is true, a problem that occupied the minds of seekers for millennia until it was rigorously reformulated by David Hilbert (1862–-1943) and definitively settled by Kurt G\" odel (1906--1978).

Hilbert's address of 1900 to the International Congress of Mathematicians in Paris is perhaps the most influential speech ever given about mathematics \cite{Hilbert1900}. In it, Hilbert outlined 23 major mathematical problems to be studied in the coming century. One of them, the 10th, raises the issue of decision in a particular case. But, in a definitive form the problem is exposed in the book that David Hilbert writes together with Wilhelm Ackermann \cite{Hilbert1928}. This book included a clear statement of the {\em Entscheidungsproblem} (decision problem): 
\begin{quote}
    {\em is there an algorithm that takes a statement formulated in a formal system and answers with "yes" or "no" whether the statement is true or false in the given system?}
\end{quote}
The first answer comes in 1931 from the logician  Kurt Gödel: 
\begin{quote}
{\bf {\em in a formal system one can correctly construct true structures whose truth cannot be proven or disproved}} \cite{Goedel1931} \cite{Goedel31}.
\end{quote}

Gödel's spectacular result triggered within five years four mathematicians Alonzo Church \cite{Church36}, Stephen Kleene \cite{Kleene36}, Emil Post \cite{Post36}, and Alan Turing \cite{Turing36, Turing1937} (listed alphabetically) to independently provide works that founded computer science. We focus in this chapter on Turing's work.

\subsection{Turing Machine}

In his paper, about "computable numbers" and {\em Entscheidungsproblem}, Alan Turing describes in words what he called a {\em computing machine}. We call it now Turing Machine (TM) and we describe it as a structure with three components (see Figure \ref{tm}a):
\begin{itemize}
    \item an infinite {\bf tape}\footnote{It is most likely that Turing was inspired by the telegraph communication systems of the 1930s that used a paper tape to receive messages.}  on/from which symbols belonging to a finite set of symbols -- the machine's alphabet -- can be write or read using
    \item an accessing write/read {\bf head} which can move in each clock cycle one position to left or right under the control of
    \item a {\bf finite automaton} (FA) whose state is updated in each clock cycle according to its state and the symbol accessed from the tape.
\end{itemize}
The translation into a representation containing digital components is shown in Figure \ref{tm}b, where the infinite tape is substituted by an infinite RAM addressed by an up-down counter. 

\placedrawing[H]{figures/12_02_turingMachine.lp}{Turing Machine. {\bf a.} Original description: Tape \& Head \& FA. {\bf b.} Digital representation: Memory \& upDownConter \& FA.}{tm}

TM is a mathematical concept because of its unlimited tape. This "infinity" is due to the fact that at the beginning of the computation the size of the space required on the tape is not known. This ignorance is dealt with by the mathematician Turing with the convenient concept of infinity. But it is beyond any doubt that any real useful computation must be done in a finite number of cycles, so we never traverse an infinite memory.


\subsection{Theoretical Foundations of Number Representations in Turing's Work}

Alan Turing’s 1937 paper, “On Computable Numbers, with an Application to the Entscheidungsproblem,” primarily focuses on the concept of computability rather than explicit number representations used in modern computer systems. However, the theoretical groundwork Turing established lays the foundation for how numbers (and more broadly, data) can be represented and manipulated by machines. Here are the key theoretical points from Turing's paper that relate to number representations:

\subsubsection{Definition of Computable Numbers}

Turing defines a \textbf{computable number} as a real number that can be calculated to any desired precision by a mechanical process. Formally, a number \( x \) is \textbf{computable} if there exists an algorithm (what Turing termed a \textbf{computing machine}, now known as a Turing Machine) that, given an integer \( n \), can compute the \( n \)-th digit of \( x \) in a finite number of steps.

In modern terms:
\[
x = \sum_{i=1}^{\infty} d_i \cdot b^{-i}
\]
where \( d_i \) are digits (either binary, decimal, or other bases) and \( b \) is the base. Turing’s framework implies that, to represent \( x \) as computable, there must be a recursive procedure for determining \( d_i \) for all \( i \).

\subsection{Representation of Numbers Using Symbols and States}

Turing’s model introduced the concept of a machine that reads, writes, and moves over a tape. The tape holds symbols from a finite alphabet, which can be used to represent numbers in various bases, including binary. Each number can thus be represented by a sequence of symbols on the tape, with the \textbf{state of the machine} directing operations on these symbols.

Formally:
\begin{itemize}
    \item Let \( \Sigma = \{0, 1\} \) represent the binary alphabet.
    \item A \textbf{string} \( w \in \Sigma^* \) on the tape can represent a binary number, where the position and value of each symbol denote its magnitude.
\end{itemize}

\subsubsection{Finite State Automaton for Number Manipulation}

Turing proposed that any computable function, including arithmetic operations on numbers, could be executed by a sequence using a finite set of states. Each pair 
\begin{verbatim} 
         {state, symbol_currently_accessed_on_tape} 
\end{verbatim} 
directs the machine to either:
1. Write or erase a symbol, 
2. Move left or right along the tape.
3. Transition to another state.

This \textbf{finite state control} is the precursor to binary arithmetic and logical operations on numbers. It implies that any number representation (binary, decimal, or otherwise) can be manipulated using a finite sequence of state transitions based on the machine's rules.

\begin{definition}
Formally, for a Turing Machine \( M(T) \), we have:
\[
M = (Q, \Sigma, \delta, q_0, F)
\]
where:
\begin{itemize}
    \item \( Q \) is a finite set of states,
    \item \( \Sigma \) is the finite tape alphabet,
    \item \( \delta: Q \times \Sigma \to Q \times \Sigma \times \{L, R, -\} \) is the transition function, where $L$ and $R$ moves the access head one position to left or right,
    \item \( q_0 \in Q \) is the initial state,
    \item \( F \subseteq Q \) is the set of final (or halting) states.
\end{itemize}
and a tape $T$ whose content is initially, in the state $q_0$,  accessed with the read/write head positioned on the first symbol.

$\diamond$
\end{definition}
This structure allows \( M(T) \) to handle number representations in binary or any base by processing symbol strings according to their state rules in a finite number of cycles. In other words, for any function $F: X \rightarrow Y$ there exists a TM whose finite automaton, starting from the initial state $q_0$ and having on the tape $x \in X$, reaches a final state, $q_f$, and generates on the tape, after a finite number of steps, $y \in Y$ if $F(x) = y$.

{\bf Important note}: for each computable sequence another TM is needed.

\subsubsection{Concept of Universal Representation}

In Turing's words:
\begin{quote}
{\em It is possible to invent a single machine which can be used to compute
any computable sequence. If this machine $\cal U$ is supplied with a tape on
the beginning of which is written the S.D of some computing machine $\cal M$
then $\cal U$ will compute the same sequence as $\cal M$.} \cite{Turing36}
\end{quote}

Turing’s concept of a \textbf{Universal Turing Machine}, UTM, suggested that any computable number or function could be represented and processed by a single machine capable of interpreting encoded instructions for any specific computation. This idea led to the realization that numbers can be encoded as \textbf{data} or \textbf{instructions}, providing a basis for multiple forms of number representation and manipulation in modern computers.

If \( p \) is a program (or representation) for a number \( x \), then \( U(p) = x \), where \( U \) is a UTM capable of simulating \( p \). Given that a finite automaton has, by definition, a complex structure, theoretical research has carefully evaluated various versions for the FA associated with the UTM. We mention two limiting cases. One is that of the work of Claude Shannon who publishes a paper describing a UTM with an automaton having only 2 states \cite{Shannon56}. The consequence of this limitation is a combinational circuit on the loop of the 2-state finite automaton of very large size. Another occurs after 50 years. It is  a paper in which the finite automaton has zero states, that is, it is no longer necessary \cite{Stefan06Turing}. This establishes the possibility of simple processors, the RISC type, processors built only from simple recursively defined circuits.  

Formally: for any computation done by a TM $M(T)$ can be used the UTM $${\cal U}(<M,T>)$$ where the stream of symbols $M$ describes the automaton of the TM $M(T)$, while the stream of symbols $T$ is the tape of the same TM. In contemporary terms, $M$ is the program, while $T$ is data , both stored in the same memory.

\subsection{Summary of Key Theoretical Points in Turing's Paper}

\begin{enumerate}
    \item \textbf{Computable Numbers}: Defined as numbers that can be calculated to any precision using a finite, mechanical process.
    \item \textbf{Finite State Control}: Established that a finite set of states could operate on symbols to represent numbers in various bases.
    \item \textbf{Machine-Based Representation}: Showed that numbers could be represented as sequences of symbols on a tape, manipulated by a machine’s states and transitions.
    \item \textbf{Universal Representation}: Introduced the idea that numbers (or their representations) could be treated as data or instructions, leading to flexible encoding schemes for number manipulation.
\end{enumerate}

Turing’s formal model set the stage for understanding how numbers can be systematically represented, stored, and manipulated by computational systems, laying foundational principles that influence number representations in modern computing.

\section{von Neumann Abstract Model}

TM is a mathematical model. The actual computing machines need an {\bf {\em intermediate model}}: an {\bf abstract} one, able to solve the problem of computation using real entities. In this "moment" comes the  serendipitous  meeting between math and tech, between UTM and digital circuits. In the abstract model proposed by the mathematician John von Neumann the finite automaton and the accessing head are integrated in what we call a processor which communicate with the memory which contains  programs and data (see Figure \ref{fig44}). 

%\section{von Neumann Computer Version: Four-Loop, Four-Order System ({\bf 4-OS})}

The two types of processors defined in the previous section were used to define two abstract computer models.

\smallskip

{\bf Historical note}: in the 1940s, the first electronic computers were made in two versions that generated two models for the subsequent evolution. Traditionally, but erroneously, these models are called {\em architectures}. The concept of architecture appeared in the 1960s. For this reason, we will call these concepts {\em abstract models} instead of architectures.

\smallskip


\placedrawing[H]{figures/05_04_von_neumann.lp}{The abstract model of computer proposed by  John von Neumann.}{fig44}

The interpretive processor is used in the definition of the {\em von Neumann abstract model}, in which the processor is loop connected to a single memory in which both programs and data are stored (see Figure \ref{fig44}). Thus, a 4-OS is obtained by connecting a 3-OS (Processor) with a 1-OS (Memory).


\section{System Organization}

The processor is the core of a computing system that contains also a memory subsystem and an input-output subsystem. The physical resources involved are tightly interleaved in various form of actual organization. The memory associated with the processor contains programs and data, while the input-output system has two main functions: expanding the memory capacity and ensuring the interaction of the system with the outside world.

\placedrawing[H]{figures/12_01_computing_components.lp}{The three-partite functionality of a computing system.}{pmio}

The way in which the tripartite functionality is distributed in a given system shows significant variations from one implementation to another. In the following, we will limit ourselves to a simple version, but not simpler than one that allows us to introduce the main concepts and problems. But, first we will show the serendipitous way in which the evolution of logic circuits meets a fundamental mathematical model for what can be rigorously defined as computation.

How looks the actual structure of a computing system? In Figure \ref{compSys} is exemplified a system where are emphasized two buses: the internal bus and the input-output bus. They are interconnected by two units: Bus Bridge and DMA (Direct Memory Access).

\placedrawing[H]{figures/06_08_computer_organization.lp}{The organization of a von Neumann version of a computing system.}{compSys}

On Internal Bus are connected fast devices while on the Input-Output Bus are connected devices involving a lot of data transfer.

The memory hierarchy is distributed in the version represented in Figure \ref{compSys}  as follows:
\begin{enumerate}
    \item Register File in the processor, having a capacity of 32W $\div$ 1024W, accessed for reading or writing in a single clock cycle
    \item Cache memory, having a capacity of 1MB $\div$ 8MB
    \item Hard Disk connected through:
        \begin{itemize}
            \item Bus Bridge
            \item DMA
        \end{itemize}
\end{enumerate}

\section{Memory Management}

%\subsection{Memory Gap}\label{memGap}

The evolution of the performance of the three components of a computer system (see Figure \ref{pmio}) occurred at a very different rate due to strictly technological aspects. The most disturbing difference, with major structural consequences, occurred in the case of processor and memory speed performances. The frequency of the processor clock increased  much faster than the memory access time (see Figure \ref{pmPGap}. Processor clock frequency increased with 60\% per year, while the access time for memories with only 10\% per year.

\placedrawing[H]{figures/06_02_processor_memory_gap.lp}{Processor-Memory Performance Gap  \cite{Patterson97}.}{pmPGap}

Due to this difference, what we now call a memory gap occurred. The fast processor cannot wait for the slow memory to deliver or receive its data. And the memory gap imposed the memory hierarchy of the computing system based on the locality principle.

\paragraph{Locality principle}\label{locprinc} says that if a certain location is accessed in the data or program memory, then the next accesses will be made with a high probability in a reasonably small vicinity.

%\bigskip

Applying the locality principle allowed the definition of the strategies that govern the conception of a hierarchy of effective memories.

\subsection{Memory Hierarchy}\label{memhier}

The hierarchy of memory is based on the fact that when we access a location that is in a memory that is too slow in relation to the speed of processing in the processor, a block or a page of bytes that contains the requested bytes will be transferred to a faster intermediate memory. This transfer is done with the hope that the following accesses will be made in the intermediate memory if the updated page or block in it was large enough.

Figure \ref{mhy} shows the current way in which the memory hierarchy is implemented. At the top of the hierarchy there are registers from register files whose content is accessible at the level of the clock signal cycle. Next comes system memory in the form of the closest level as cache memory. The name comes from the French {\em cach\' e} which means hidden. Indeed, this level of memory/memories is transparent to the architecture of the computing system in most cases.

\placedrawing[H]{figures/06_03_memory_hierarchy.lp}{Memory hierarchy from fast and small register files to slow and large auxiliary storage memories such as SSD (solid-state drive), HDD (hard disk drive),  optic disk, magnetic tape.}{mhy}

Registers  are managed by compilers, their size is $< 4KB$ (in most of cases $\sim 128B$, i.e., 32 32-bit words), the access time is $< 0.5 ns$.

Cache memories are managed by hardware, their size is $< 16MB$, the maximum access time is $\sim 0.5 ns$.

The main memory is managed by the operating system, its size is in the range of $1 \div 16 GB$, the access time is $\sim 100 ns$.

The disk memory is managed by the operating system or human operator, its size is $> 128GB$, the access time is $\sim 5 ms$.

\subsection{Virtual Memory Mechanism}

How relate two successive levels in the hierarchy? Let us consider two levels in the memory hierarchy, $Level_i$ containing $m$ pages and $Level_{i+1}$ containing $n$ pages, with $n >> m$. Each page in $Level_i$, called {\em physical level} is a copy of a page from $Level_{i+1}$ called {\em virtual level}. In Figure \ref{vmm}a, pages 0, 2, and $m-2$ from the level $i$ were brought from pages 3, 4, and 0 of level $i+1$. Any change in the physical level must actualized in the associated virtual level. The "user"  accesses the virtual level while the virtual mechanism accesses the physical level. It is the job of the virtual memory mechanism to perform transparently the access and to keep coherent the content of the virtual level with the physical level.

\placedrawing[H]{figures/06_04_virtual_memory.lp}{Virtual memory mechanism. {\bf a.} Page correspondence. {\bf b.} Translation mechanism.  }{vmm}

The addressing mechanism is represented in Figure \ref{vmm}b, where the access to the physical memory is done mediated by a translation mechanism performed by the {\em Page translate} block.  The "user" issues a virtual address having two parts:
\begin{itemize}
    \item virtual page address
    \item offset in the page
\end{itemize}
When a new page is loaded into physical memory, the content of this translator is updated according to the correspondence shown in the figure \ref{vmm}a.

Page translation mechanism is implemented in various form.

The simplest one is a RAM memory containing $n$ $log_2 m$-bit words. The RAM memory is in this case almost empty, because only $m$ locations from  $n$ are used, because usually $n>>m$.

A more efficient solution is to use an associative memory to make the translation.

\subsection{Associative Memory-Based Lookaside Page Translator (TLB)}

The associative memory (see \ref{assocMemory} in this book)  allows us to design a translation lookaside buffer (TLB) with a size given by the number of pages in the physical memory, unlike the one based on a RAM type memory which has a size given by the much larger number of pages in the virtual memory.

\subsubsection{RAM-Based Translation Lookaside Buffer (TLB)}

One simple solution for {\em Page translate} block in Figure \ref{vmm}b is to use a RAM with $(log_2 m)$-bit $n$ locations. In each location can be stored the address of a physical page.  Because $n >> m$ almost all locations in this RAM contain in each moment a meaningless information.In the example shown in figure \ref{vmm}a, $RAM(0) = m-2$, $RAM(3) = 0$, $RAM(4) = 2$, $\ldots$, so the RAM usage as TLB is at most $m/n\%$. 

\subsubsection{Associative Memory-Based Translation Lookaside Buffer (TLB)}

A less simple solution than using a RAM is to use an associative memory (AM). The dimensioning of this version requires a solution in which the entire physical structure is usable. The storage array has $m$ locations of $log_2 m$ bits each and the {\em programmable decoder} (the CAM) also has $m$ locations of $log-2 n$ bits each. So the entire structure has a dimension in $O(m \times (log \; m + log \; n)$ compared to solutions using RAM which impose a dimension in $O(n \times (log \; m + log \; n))$.

The gain in size, i.e., the area occupied by the two solutions, with RAM or with AM, is:
$$\frac{S_{TLB}(RAM)}{S_{TLB}(AM)}  \in O(n/m)$$
Even though a content-addressable memory cell (see Figure \ref{cam_cell}) is more elaborate and larger, because the ratio $n/m$ is sufficiently large, the AM solution, which uses a CAM as a programmable decoder, is preferable.

It should not be neglected that the smaller size of the associative memory allows for an easier achievement of the speed performance of the structure.

\subsubsection{Translation Lookaside Buffer (TLB)}

With a minimal number of differences, the described virtualization mechanism is applied in managing the relationship between all the hierarchical levels in the memory of the computer system.

In what follows, we will consider the simple example in which the hierarchy contains a single cache level, the main memory and the hard disk (see Figure \ref{compSys}). We will compare the main parameters that characterize the  cache-main memory relationship and the main memory-hard disk relationship\footnote{In what follows, the numerical values ​​evolve very dynamically over time. Which is why they must be considered only through the relative aspects they highlight.}.
\begin{itemize}
    \item the typical amount of data moved when physical memory is actualized from the virtual memory
        \begin{itemize}
            \item block size for cache: 128B
            \item page size main memory: 4KB
        \end{itemize}
    \item the typical access time when the addressing is hit
        \begin{itemize}
            \item 1 clock cycle (very rarely 2 or 3) for cache
            \item $\sim 100$ clock cycles because of board technology and silicon technology for memories (the first is reduced by multi-chip packaging) for main memory
        \end{itemize}
    \item the typical miss penalty:
        \begin{itemize}
            \item 100 clock cycles for cache memory
            \item $10^6$ clock cycles for main memory
        \end{itemize}
    \item the typical miss rate:
        \begin{itemize}
            \item $1\%$ for cache memory
            \item $10^{-4}\%$ for main memory
        \end{itemize}
\end{itemize}



\section{I/O}

The input-output system includes certain specific mechanisms that we will briefly describe first. Then we will briefly describe the most important input-output devices.

\subsection{Bus: the Interconnection Fabric}

A {\bf {\em Bus}} is a communication system that transfers data between components inside a computer, or between computers. There are two types of buses: serial (those that transfer data bib by bit) and parallel (those that transfer data word by word of n bits). The main problem that arises is that of the synchronization of the transfer. We will discuss it under its essential aspects for the synchronous transfer of n-bit words.

\placedrawing[H]{figures/06_09_read_write_waveform.lp}{The waveforms for a read cycle followed by a write cycle.}{bus}

In Figure \ref{bus}, the waveforms for data transfer on the Internal Bus (see Figure \ref{compSys}) between Main Memory and 2nd Level Cache are presented. Three important moments are marked in the mentioned waveforms:
\begin{description}
    \item[$t_1$]: the active edge of clock takes a read command (Read = 0) from the address Read Address. Both Read and Read Address must be stable at least minimum {\bf {\em  set-up time}} before the active edge of clock (the positive one) and {\bf {\em hold time}} after the positive edge of clock
    \item[$t_2$]: the active edge of clock can take the data from the output of Main Memory only if the signal Wait is not active (Wait = 0) thus completing the read cycle. The signal Wait is necessary when the Main Memory is not a fast enough memory, or the clock frequency is too high.
    \item[$t_3$]: the active edge of clock takes Valid Input Data to write it at Write Address because Write signal is active (Write' = 0).
\end{description}


The Wait signal is deactivated (Wait = 0) at $t_2$ introducing a latency of 2 clock cycles in the example illustrated by Figure \ref{bus}. Usually, this latency is very high, reaching hundreds of cycles.



\subsection{DMA}

\subsubsection{Definition}

Direct memory access (DMA) (see Figure \ref{compSys}) is an important feature of computer systems. It allows  to access the main memory  system independently of the central processor.

When the processor is running an input/output program, it is occupied for the entire duration of the read or write operation. Thus the processor is unavailable to perform other work. With DMA, the processor  initiates the transfer only, then it does other operations while the transfer is performed. When the transfer is completed the processor receives an interrupt from the DMA (about interrupt see \ref{interrupt} in this lecture notes).

An example of using the DMA unit is for transferring a page of 4KB from Hard Dist to Main Memory.

The CPU initializes the DMA by sending the given information through the data bus:
\begin{itemize}
    \item the starting address of the memory block where to be read or write.
    \item the number of words in the memory block
    \item control bit to define the mode of transfer, read or write.
    \item control bit to begin the DMA transfer
\end{itemize}

The DMA controller transfers the data in three modes:
\begin{itemize}
    \item burst mode:
    \begin{itemize}
        \item entire data or burst of block containing data is transferred before CPU takes control of the buses back from DMA
        \item is the quickest mode of DMA Transfer since at once a huge amount of data is being transferred
        \item since at once only the huge amount of data is being transferred so time will be saved in huge amount
        \item {\bf pros}: fastest mode of DMA Transfer
        \item {\bf cons}: less user friendly because during the DMA transfer CPU will be blocked.
    \end{itemize}
    \item cycle stealing mode for each transfer:
    \begin{itemize}
        \item slows I/O device because will take some time to prepare data  and within that time CPU keeps the control of the buses.
        \item once the data is ready CPU give back control of system buses to DMA for 1-cycle in which the prepared word is transferred to/from memory.
        \item as compared to burst mode this mode is little bit slowest since it requires little bit of time which is actually consumed by I/O device while preparing the data.
        \item {\bf pros}: most Efficient way for DMA Transfer.
        \item {\bf cons}: CPU won’t be blocked entire time.
    \end{itemize}
    \item transparent or interleaving mode:
    \begin{itemize}
        \item whenever CPU does not require the system buses, then the control of buses will be given to DMA
        \item CPU will not be blocked due to DMA at all
        \item it is the slowest mode of DMA transfer since DMAC has to wait might be for so long time to just even get the access of system buses from the CPU itself
        \item less amount of data will be transferred
        \item {\bf pros}: CPU will not be blocked at all
        \item {\bf cons}: Slowest DMA transfer rate.
    \end{itemize}
\end{itemize}

\placedrawing[h]{figures/06_12_init_trans.lp}{Initiating the transfer between I/O and Memory.}{tranferInit}

The transfer between an I/O device and MEMORY under the control of the DMA is initiated by the next sequence mode:
\begin{enumerate}
    \item I/O sends request DMA, {\bf DRQ}, and DMA sends to CPU the hold request, {\bf HLD}
    \item CPU acknowledge to DMA with {\bf HLDA}
    \item DMA acknowledge to I/O with {\bf DACK}
\end{enumerate}
and when data transfer is done,  DMA sends an {\em interrupt}, INT, to PROCESSOR (see Figure \ref{tranferInit}).


\subsubsection{Structure}

The structure of a minimal DMA connected between two buses, BusA and BusB, contains three main modules:
\begin{itemize}
        \item a buffer for the data being transferred sometimes between BusA and  BusB and other times between BusB and  BusA
        \item a presettable counter for the memory address from which or to which data is transferred
        \item a counter extended automaton to control the transfer process.
\end{itemize}
The resulting block schematic is represented in Figure \ref{simpleDMA}, where:

\placedrawing[h]{figures/06_13_minimal_DMA.lp}{Block schematic for a minimal DMA.}{simpleDMA}

\begin{itemize}
    \item CEA is the counter extended automaton whose counter is initiated with data from BusA under the Processor command; the automaton of CEA recives and sends signal from/on both buses and generates commands for the multiplexer and FIFO
    \item FIFO is used as a buffer to optimize the transfer time between the two busses where the speed of transfer are very different; the block of data on BusA is transferred faster then the same block of date on BusB because Memory is much faster then the I/O device
    \item the multiplexor, mux, allows to use the same FIFO for the transfers in both directions, from I/O to Memory or from Memory to I/O
    \item AddressCounter is used to generate the addresses for Memory; it is loaded by Processor in the initialization process of a transter
\end{itemize}
Representation from Figure \ref{simpleDMA} is a functional schematics and does not contains the interface circuits between DMA and the two buses.

The critical timing in the transfer is on the bus BusA where DMA compete with Processor in using the bus for communicating with Memory in the computing process. For this reason, the transfer mode on BusA is burst mode facilitated by the  FIFO memory. When I/O transfers a bloc in Memory, first the FIFO memory is loaded with a block of data at the low speed of the I/O device, then the FIFO memory is emptied in a burst performed at the high speed of the Memory module. Symmetrical, the FIFO memory is fast loaded from Memory in burst mode, then the data is transferred  at the low speed imposed by the I/O device on the BusB. The overall time is imposed by the slowest device, the I/O system, but the time during which the processor is deprived of the possibility of using the BusA bus is minimized by the appropriate use of the FIFO memory.

In addition to the signals (read, write, sel, load, increment) and flags (empty, full) that the CEA generates inside the DMA, the control signals (commands and flags) on the two buses are also managed by the same circuit, as follow:
\begin{itemize}
    \item on BusA:
        \begin{itemize}
            \item commands: {\tt HLD, INT, readA, writeA, startReadA, startWriteA}
            \item flags: {\tt HLDA, waitA}
        \end{itemize}
    \item on BusB:
        \begin{itemize}
            \item commands: {\tt DACK, readB, writeB, startReadB, startWriteB}
            \item flags: {\tt DRQ, waitB}
        \end{itemize}
\end{itemize}
In addition to these signals, the CEA finite state machine also has signals related to the associated presettable counter. These are: the {\tt zero} flag (which indicates that the counter has reached zero by decrementing), and the {\tt preset} (loads the transfer size minus 1 in presetable counter) and {\tt dec} (decrements the counter value on each write or read).

\begin{exemplu}\label{minDMA}
The process of transferring a block of data from or to a I/O device using a minimal DMA structure is presented. It is considered that the system works with two buses, an internal bus, BusA, and an input-output bus, BusB. On the internal bus we need to protect the computing process from too intensive use of the bus by DMA. For this we will use the BusAt transfer in burst mode. This type of transfer negotiates the takeover of the bus control only once, and the transparent mode increases the transfer time too much.

{\em
{\small
\begin{shaded*}
\begin{lstlisting}[language=Verilog, caption= , morekeywords={always_ff, always_comb, logic, loop, repeat, doInParallel, for}]
/************************************************************************
The transfer procedure has only two modes: (1) read from I/O to memory
a block of maximum 4096 words or (2) write from memory to I/O a block
of maximum 4096 words.
************************************************************************/
state0: idle state, waiting for starting a read/write transfer
        {nextState,com} = startReadA  ? {state1,load} :
                         (startWriteA ? {state2,load} : (state0,preset))

state1: Processor loads in FIFO the source block address in I/O
        {nextState,com} = {state3,write}

state3: Processor loads in AddressCounter the starting address in Memory
        {nextState,com} = {state4,load}

state4: I/O receives from FIFO the pointer to the block
        (nextState,com} = {state5,read,writeB}

state5: the block to be transferred is loaded iin FIFO
        {nextState,com} = waitB ? {state5,-}                       :
                         (zero  ? {state11,readB,write,dec,preset} :
                                  {state5,readB,write,dec})

state11: DMA request control on BusA
        {nextState,com} = HLDA ? {state6,HLD} : {state11,HLD}

state6: the bloc is transferred in burst mode from FIFO to Memory
        {com,nextState} = {writeA,(waitA ? state6 :
                                  (zero ? {read,incAddr,dec,state0} :
                                          {read,incAddr,dec,state6}))}

state2: Processor loads in FIFO the destination block address in I/O
        {nextState,com} = {state7,write}

state7: Processor loads in AddressCounter the starting address in Memory
        {nextState,com} = {state8,load}

state8: I/O receives from FIFO the pointer to the block
        (nextState,com} = {state9,read,writeB}

state9: the block is transferred in burst mode from Memory to FIFO
        {com,nextState} = {readA,(waitA ? state9 :
                                  (zero ? {write,incAddr,dec,state10} :
                                          {write,incAddr,dec,state9}))}

state10: the block is transferred from FIFO to I/O
        {nextState,com} = waitB ? {state10,-}                      :
                         (zero  ? {state0,writeB,write,dec,preset} :
                                  {state10,writeB,write,dec})
\end{lstlisting}
\end{shaded*}
}}

$\diamond$
\end{exemplu}



\subsection{I/O Devices}

There are two types of I/O devices:
\begin{itemize}
    \item Storage devices, mainly represented by hard disk, flash memory, optical disk, tape.
    \item Functional devices, mainly represented by display system, network interface, ...
\end{itemize}
We will briefly present the storage devices that extend the memory hierarchy outside the central computing unit.

\subsubsection{Hard Disk}

The most common auxiliary memory is the hard disk drive. Its internal structure is represented in Figure \ref{disk}.

\placedrawing[H]{figures/06_11_hard_drive.lp}{Hard Disk Drive Structure [Source: Google-Images]}{disk}

Main parameters:
\begin{itemize}
    \item storage capacity: 10 $\div$ 30 TB
    \item average access time: 2.5 $\div$ 10 ms
    \item MTBF: $\sim$285 years (MTBF stands for mean time between failures; it is the predicted elapsed time between inherent failures of a mechanical or electronic system, during normal system operation.)
    \item bytes per sector: 512 $\div$ 4096
    \item shock tolerance
        \begin{itemize}
            \item operating: 10 $\div$ 100 g (g stands for gravitational acceleration: $9.8 m/s^2$.)
            \item nonoperating: 100 $\div$ 100 g
        \end{itemize}
\end{itemize}

\subsubsection{Optical Disks}

Optical disks are read-only mediums. There are two versions:
\begin{itemize}
    \item CD: optical compact disk (0.65 GB)
    \item DVD: digital video disk (4.7 GB); sometimes written on both sides to double capacity
\end{itemize}

\subsubsection{Flash Memory}


Flash memory is a type of EEPROM (electronically erasable programmable readonly memory), which is normally read-only but can be erased. Solid-state devices (SSDs) based on flash memory technology do not have moving parts.

The capacity per flash memory chip increased by about 50\%–60\% per year.

Typical SSDs will have a the time taken for a disk drive to locate the area on the disk where the data to be read is stored between 0.08 and 0.16 ms.

Time to {\bf read} 2 KB from a flash memory takes about 75 $\mu$S,  while DDR SDRAM takes less than 500 ns. Then  flash memory is  about 150 times slower. But compared to HDD is $\sim 400$  times faster. We conclude: the flash memory can not replace DRAM for main memory, but is a good candidate to replace HDD.

Time to {\bf write} for flash memory is about 1500 times slower then SDRAM, and about 8–15 times faster than HDD.

\subsubsection{Tapes}

Magnetic-tape data storage is a system for storing  information on magnetic tape using digital recording.

Typically  9-track tape are used to store bytes with CRC. Magnetic tape  packaged in cartridges and cassettes. The device that performs the writing or reading of data is called a tape drive.

Tape data storages are now used mainly for system backup, data archive or data exchange.

Uncompressed/Native capacity: $\sim$20TB.

Compressed capacity: $\sim$50TB.




